% ch-a5-03.tex —— 第 5 篇（Ben-Zvi--Chen--Helm--Nadler）续片，底本 PDF 458--469 = 印刷页 445--456
\ysrc{445}

\item 对 $G$-等晶体（isocrystal）叠上 $\ell$-进层的分次形式（[HL22] 意义下）构造一个 $S^1$-作用。
\item 把 Fargues 与 Zhu 诸猜想的混合形式提升为 $S^1$-等变形式，即给出这些分次范畴的一个 $S^1$-等变的等同——参见 [Zhu25] 中的注 4.6.8。
\item 将自守侧可 $S^1$-等变化（$S^1$-equivariantizable）的层构成的满子范畴，与对应于纯内型表示的子范畴作比较。
\item 说明：该子范畴上凝聚局部 Langlands 对应的循环形变，在固定无穷小特征时特化为 Vogan 可构造局部 Langlands 对应的（范畴形式）。
\end{itemize}
\end{prob}

\begin{rem}{3.10}\label{rem:a5-03-3-10}
人们或许会想忘掉 $\GG_m$-等变性，改成只对那个自同构要求 $\Z$-等变性，以此一般地引入自同构迹上的 $S^1$-作用。具体地说，设 $X$ 是带自同构 $\phi$ 的叠；可以构造叠 $X/\phi^{\Z}$，它是 $B\Z\simeq S^1$ 上的叠。于是其圈空间 $\mathcal L(X/\phi^{\Z})$ 位于 $\mathcal L(B\Z)\simeq \Z\times S^1$ 之上。取 $\{1\}\times S^1$ 上的纤维（换言之，取 $\Map(S^1,X/\phi^{\Z})$ 中沿 $X/\phi^{\Z}\to B\Z=S^1$ 的复合为次数 $1$ 的映射构成的子叠），显然得到带 $S^1$-作用的空间。遗憾的是，所得 $S^1$-作用自由，于是所得范畴全为 $u$-挠的，从而被 $S^1$-局部化消灭。粗略地说，$\Z$-等变性与 $\GG_m$-等变性之别，可从函子 $\Rep(\GG_m)\to\Rep(\Z)$（即沿把 $1$ 映到所取自同构的映射 $\Z\to\GG_m$ 的拉回）不再全忠实看出：后者有非平凡的 $\mathrm{Ext}^1$，而前者没有。
\end{rem}

\section{从阿基米德局部 Langlands 到扭量几何 Langlands}

本节讨论如何在阿基米德情形实现本文的基本范式：借由导出圈空间上圆作用的等变局部化，把局部 Langlands 对应的凝聚形式与可构造形式联系起来（沿 [BN07, BN10b, BN16]）。为此须先描述对应的这两种形式。

我们所取的可构造形式是 Soergel 猜想的一个变体；凝聚形式则由扭量线 $\ol\PP^1_\R$ 上驯顺分歧的几何 Langlands 对应给出，后者可视为 Fargues 猜想的阿基米德对应物。

我们从较熟悉的、可构造的阿基米德局部 Langlands 对应（固定无穷小特征）说起。它在 Grothendieck 群层面是 Adams--Barbasch--Vogan（ABV）的定理，在导出范畴层面则是 Soergel 的猜想。第 4.1 节描述自守侧：用 Kashiwara--Schmid 版的 Beilinson--Bernstein 局部化，把纯内型的相关表示范畴实现为旗簇上等变可构造层的范畴。第 4.2 节（沿 [BN10b]）给出谱侧即 ABV 几何参数空间 [ABV92] 的一个简洁的叠式描述（另见 Mason-Brown 在本文集中的文章 [MB25, 2.2]）。两者的比较即 Soergel 猜想，见第 4.3 节。

\ysrc{446}

接着我们陈述阿基米德局部 Langlands 的「族」版本，其中允许无穷小特征变动。第 4.4 节描述自守（表示论）侧，第 4.5 节描述谱侧。具体地说，我们引入 Langlands 参数的叠 $\mathbb L^\eta$，它是非阿基米德情形所遇诸叠的阿基米德对应物。这个叠有一个初等而显式的描述，令人联想到幂幺 Langlands 参数的叠；但它也可实现为扭量线 $\ol\PP^1_\R$ 上在无穷远驯顺分歧的局部系统的叠，这使它作为扭圈空间自然带有圆作用。我们说明：借由圈的一种 Jordan 分解，这个叠在变动的无穷小特征所对应的 ABV 参数空间之间光滑地插值。然后我们陈述 Soergel 猜想的「族版本」：它把纯内型（无穷小特征任意）的光滑表示等同于 $\mathbb L^\eta$ 上的循环层（即 $S^1$-等变凝聚层上的 Tate 构造）。这个猜想与上一章所述的非阿基米德图景密切平行。

最后，第 4.6 节解释 $\mathbb L^\eta$ 上凝聚层全范畴的自守对应物。我们（沿 [BN07, BN16]）陈述 $\ol\PP^1_\R$ 上驯顺分歧的几何 Langlands 猜想，并说明它如何把 Soergel 猜想作为其周期循环形变而恢复出来。自守侧背后的几何机制是：把 Soergel 猜想的自守侧出现的那些等变旗簇（经由表示的 Kashiwara--Schmid 描述）实现为 $\ol\PP^1_\R$ 上抛物丛叠的半稳定轨迹（以及 $S^1$-不动点）。

\subsection{阿基米德局部 Langlands：自守侧}

设 $(G,\theta)$ 是带拟分裂实形式的复约化群。实形式 $\theta$ 给出纯内型的一个集合 $\Theta$，它们在几何上如下出现 [BN10b, Ber16]：考虑 $BG$ 的伽罗瓦不动点 $BG^\Gamma$，其中 $\Gamma$ 经由共轭 $\theta$ 作用。则有分解
\[
BG^\Gamma\simeq\bigsqcup_{\tau\in\Theta}BG_\tau,
\]
其中 $\Theta$ 是 $\theta$ 的纯内型的等价类之集，$G_\tau$ 是相应的实形式（可以带重数出现）。

局部 Langlands 对应为群 $G_\tau$（$\tau$ 取遍 $\Theta$）参数化表示。因此实局部 Langlands 纲领的终极目标之一，是用 Langlands 对偶的语言描述纯内类 $\Theta$ 中实群的 Harish-Chandra 模的整个 dg 范畴
\[
\mathcal{HC}_\Theta=\bigoplus_{\tau\in\Theta}\mathcal{HC}_\tau.
\]
对每个 $[\lambda]\in\frakh^*/W$，以 $\mathcal{HC}_{\tau,[\lambda]}$ 记实形式 $G_\tau$ 的、带 pro-完备化广义无穷小特征 $[\lambda]$ 的 Harish-Chandra 模的 dg 范畴，并令
\[
\mathcal{HC}_{\Theta,[\lambda]}=\bigoplus_{\tau\in\Theta}\mathcal{HC}_{\tau,[\lambda]}.
\]
现在设 $\lambda$ 正则（奇异情形的讨论见注 4.3），并令 $K_\tau$ 记 $G_\tau$ 的极大紧子群的复化。于是可以应用 Beilinson--Bernstein 局部化，用旗簇 $G/B$ 上对相应复对称子群 $K_\tau$ 等变的 $\lambda$-扭 D-模来描述实群 $G_\tau$ 的 Harish-Chandra $(\frakg,K_\tau)$-模范畴 $\mathcal{HC}_{\tau,[\lambda]}$：
\ysrc{447}
\[
\mathcal{HC}_{\tau,[\lambda]}\simeq \mathcal D_\lambda(K_\tau\backslash G/B).
\]
由 Riemann--Hilbert 对应，这个 de Rham 实现有 Betti 对应物：$K_\tau\backslash G/B$ 上 $\alpha=\exp(2\pi i\lambda)$-扭的可构造层（见 [K89]）。

Kashiwara 与 Schmid [KS94, K08] 引入的另一等同，把表示的导出范畴与旗簇上的等变导出范畴联系起来，它与群 $G_\tau$ 的容许表示关系更直接（例如通过自然的整体化函子）：
\[
\mathcal{HC}_{\tau,[\lambda]}\simeq \mathrm{Shv}_\lambda(G_\tau\backslash G/B).
\]
在这个实现中，我们把 $K_\tau$-等变性换成了 $G_\tau$-等变性；两种实现的等同由层的 Matsuki 对应 [MUV92] 给出。

借由同伦不动点，可以同时描述变动 $\tau$ 时的所有这些范畴：有自然的等同
\[
\bigsqcup_{\tau\in\Theta}G_\tau\backslash G/B\simeq (B\backslash G/B)^\Gamma,
\]
其中 $\Gamma$ 先交换两个因子、再复合实形式 $\theta$，以此作用在 $B\backslash G/B\simeq G\backslash(G/B\times G/B)$ 上。换言之，整个纯内类 $\Theta$ 的表示被自然地实现为
\[
\mathcal{HC}_{\Theta,[\lambda]}\simeq\mathrm{Shv}_\lambda((B\backslash G/B)^\Gamma).
\]

\subsection{阿基米德局部 Langlands：谱侧}

设 $(\check G,\eta)$ 是 $(G,\theta)$ 的 Langlands 对偶群，带其对偶代数对合。经由伽罗瓦等变 Satake 等价 [Zh15, Zh17]，引入与 $G$ 及共轭 $\theta$ 相伴的几何 L-群 $G^L$。它是 $\check G$ 被伽罗瓦群 $\Gamma=\Gal(\C/\R)$ 扩张所得的扩张 $\check G\to G^L\to\Gamma$（但未必是半直积）。为避免 L-群定义中的细微之处，我们假设 $G$ 为伴随型。

在谱侧，对合 $\eta$ 同样给出 $\eta$ 的内类中一族对合，由集合 $\Sigma(\eta):=\{\sigma\in\check G\mid\sigma\eta(\sigma)=e\}/\check G$ 参数化；同样有分解
\[
B\check G^\Gamma\simeq\Hom(\Gamma,\check G\rtimes\Gamma)/(\check G\rtimes\Gamma)\simeq\{\sigma\in\check G\mid \sigma\eta(\sigma)=e\}/\check G\simeq\bigsqcup_{\sigma\in\Sigma(\eta)}BK_\sigma,
\]
其中 $K_\sigma:=\check G^\iota\subset\check G$ 是对合 $\iota(g)=\tilde\sigma\eta(g)\tilde\sigma^{-1}$ 的对称子群，该对合由代表元 $[\tilde\sigma]=\sigma\in\Sigma(\eta)$ 给出。

以 $X(G^L)$ 记几何参数的 ABV 空间。ABV 局部 Langlands 对应只涉及 $X(G^L)$ 上的等变层，故自然应转而考虑商叠 $\mathcal X(G^L):=X(G^L)/\check G$。叠 $\mathcal X(G^L)$ 是诸叠 $\mathcal X(\mathbb O_\lambda,G^L)$ 在诸半单轨道 $\mathbb O_\lambda=\check G\cdot\lambda\subset\check{\frakg}$ 上的不交并。记 $e(\mathbb O_\lambda)=\exp(2\pi i\mathbb O_\lambda)=\check G\cdot\alpha\subset\check G$，其中 $\alpha=\exp(2\pi i\lambda)$；并记 $\check P(\lambda)\subset\check G(\alpha)$ 为中央化子中相应的抛物子群。我们记对称子群 $K_\sigma(\alpha):=P(\lambda)^\sigma$\tn{底本此处作 $P(\lambda)$（不加 check），而第 449 页（注 4.3 中）同一对象记作 $\check P(\lambda)$；两处记号不一致，疑有一处漏写 check。译文照排。（第 447 页）}，其中
\begin{equation}\tag{4.1}\label{eq:a5-03-4-1}
\sigma\in\Sigma(\eta,\alpha):=\{\sigma\in\check G\mid\sigma\eta(\sigma)=\alpha\}/\check G(\alpha).
\end{equation}

\ysrc{448}

关于 $\check G$-轨道的描述 [MB25] 与同伦不动点的一些简单观察相结合，给出 [BN10b] 中等变 ABV 空间的如下简洁描述。

\begin{prop}{4.1}\label{prop:a5-03-4-1}
对固定的 $\lambda$，ABV 叠 $\mathcal X(\mathbb O_\lambda,G^L)$ 等同于不交并
\[
\bigsqcup_{\sigma\in\Sigma(\eta,\alpha)}K_\sigma(\alpha)\backslash\check G(\alpha)/\check P(\lambda),
\]
而后者又等同于不动点叠 $(\check P(\lambda)\backslash\check G(\alpha)/\check P(\lambda))^\Gamma$，其中 $\Gamma$ 交换两个因子并复合 $\eta$ 而作用。
\end{prop}

\subsection{Soergel 猜想}

现在设 $\mathbb O$ 是正则半单轨道，即我们参数化的是带正则无穷小特征的表示；奇异特征的讨论见注 4.3。特别地，$\alpha$ 的中央化子 $\check G(\alpha)$ 中相应的抛物子群 $P(\lambda)$ 是 Borel 子群 $\check B(\alpha)$。此时，ABV 参数化可以简单地表述为：描述 $(\check B(\alpha)\backslash\check G(\alpha)/\check B(\alpha))^\Gamma$ 上层范畴的 Grothendieck 群。

\begin{conj}{4.2}（Soergel 猜想，[BN10b] 的表述）\label{conj:a5-03-4-2}
存在 Koszul 对偶
\[
\mathrm{Shv}_\alpha((B\backslash G/B)^\Gamma)\underset{Kos}{\sim}\mathrm{Shv}((\check B(\alpha)\backslash\check G(\alpha)/\check B(\alpha))^\Gamma),
\]
它联系纯内类 $\Theta$ 中、带正则无穷小特征 $\lambda$（指数为 $\alpha$）的 Harish-Chandra 模的范畴与 ABV 几何参数叠上的层范畴。
\end{conj}

所谓 Koszul 对偶 $\mathcal C\underset{Kos}{\sim}\mathcal D$，是说范畴 $\mathcal C$ 与 $\mathcal D$ 有分次提升 $\mathcal C_{gr}$、$\mathcal D_{gr}$（亦称「混合版本」），且在做一次上同调剪切（cohomological shear，见定义 A.10）之后，有分次范畴的等价
\[
\mathcal C_{gr}\simeq\mathcal D^{/}_{gr}.
\]
因为有了这次剪切，底层的范畴 $\mathcal C$、$\mathcal D$ 未必等价；但它们的 2-周期版本（周期局部化）是等同的：
\[
\mathcal C\otimes_k k[u,u^{-1}]\simeq\mathcal D\otimes_k k[u,u^{-1}].
\]
这样的 Koszul 对偶诱导 $\mathcal C$ 与 $\mathcal D$ 的 Grothendieck 群之间的对偶（例如见 [BV21] 的引言），于是 Soergel 猜想给出 Vogan 的特征对偶 [Vo82]。

猜想 4.2 的一个更强形式要求：该等价尊重两侧 Hecke 范畴的自然作用，从而把表示 Grothendieck 群的 Hecke 代数对称性范畴化（「凝聚延拓表示」）。事实上，当 $\alpha=1$ 时，[BY13] 的 Koszul 对偶定理给出幺半等价
\ysrc{449}
\[
\mathrm{Shv}_1(B\backslash G/B)\underset{Kos}{\sim}\mathrm{Shv}(\check B\backslash\check G/\check B),
\]
它联系 $N\backslash G/N$ 上的幂幺双单值层（unipotent bi-monodromic sheaves）与对偶旗簇上的等变层——即 [BGS96] 的 Koszul 对偶的强幺半形式，也就是 Soergel 猜想在复情形、平凡单值下的情形（已在 [So01] 中证明）。较近的结果是 [LuY20] 的内窥 Koszul 对偶定理，它给出幺半等价\footnote{对中心连通的 $G$（如伴随型）成立；否则左边须限制在主块上。}
\[
\mathrm{Shv}_\alpha(B\backslash G/B)\underset{Kos}{\sim}\mathrm{Shv}(\check B(\alpha)\backslash\check G(\alpha)/\check B(\alpha)),
\]
从而给出 Soergel 猜想一般（正则）复情形的幺半形式。猜想 4.2 的两侧都是此处等同的两个幺半范畴的自然模范畴，且该等价理应尊重这一结构。

\begin{rem}{4.3}（奇异无穷小特征）\label{rem:a5-03-4-3}
当 $\lambda$ 奇异时，须对猜想 4.2 的陈述作如下修改。在谱侧，把 Borel 子群 $\check B(\alpha)$ 换成由 $\lambda$ 决定的抛物子群 $\check P(\lambda)$（如 [MB25] 所述）。在自守侧，$\lambda$ 奇异时 Harish-Chandra 模的范畴是旗簇上相应的等变层范畴的商，其描述例如见 [K89]。等价地，有一个自然投射子（幂等单子）作用在层范畴上，其模恰是表示。在 Hecke 范畴的层面（或在复群情形），这些互相匹配的修改正是 [BGS96, BY13] 的抛物-奇异 Koszul 对偶。
\end{rem}

\subsection{实群的表示族}

现在我们重新表述 Soergel 猜想的自守侧，即实形式 $G_\tau$（落在 $\theta$ 的纯内类中、带正则无穷小特征 $\lambda$）的 Harish-Chandra 模的范畴 $\mathcal{HC}_{\Theta,[\lambda]}$。如前述，Kashiwara--Schmid 局部化 [KS94, K08] 把它实现为 $\alpha$-扭的层（$\alpha=\exp(2\pi i\lambda)$），所在空间为
\[
\bigsqcup_{\tau\in\Theta}G_\tau\backslash G/B\simeq (B\backslash G/B)^\Gamma.
\]
为了容纳 $\alpha$ 的变动，我们先把 $G/B$ 上 $\alpha$-扭的层改写为环面丛 $G/N\to G/B$ 上、沿纤维局部常值且带广义单值 $\alpha$ 的层。要让 $\alpha$ 变动，只需去掉单值性条件，允许沿纤维局部常值的任意层。这样的局部常值性可以改述为：不允许层在奇异支撑中出现半单方向——也就是说，我们考虑的是 $\mathrm{Shv}_N(G/N)$，即 $G/N$ 上带幂零奇异支撑的层。

于是，猜想 4.2 自守侧的一个族版本，由对纯内类 $\tau\in\Theta$ 取直和所得的诸幂零层范畴 $\mathrm{Shv}_N(G_\tau\backslash G/N)$ 给出。

\subsubsection{从层到表示}

自然会问：范畴 $\mathrm{Shv}_N(G_\tau\backslash G/N)$ 与我们的初衷——即实群 $G_\tau$ 的表示——有何关系？为此需要某种中心特征变动的 Kashiwara--Schmid 局部化。

\ysrc{450}

在 D-模的 de Rham 场合，[BN12] 引入了无穷小特征变动的 Beilinson--Bernstein 局部化版本：把 $G/B$ 上的扭 D-模 $\mathcal D_\alpha(G/B)$ 换成 $G/N$ 上弱 H-等变的 D-模 $\mathcal D_H(G/N)$。结果是等同
\[
U\frakg\text{-}\mathrm{mod}\simeq\mathcal D_H(G/N)^{\mathbb D},
\]
它把 $U\frakg$-模等同于 Weyl--Demazure 单子 $\mathbb D$ 的模；$\mathbb D$ 是作用在 $\mathcal D_H(G/N)$ 上、强制 Weyl 群不变性的一个显式代数——也就是说，它负责从依赖 $\lambda\in\frakh^*$ 的 D-模到依赖 $[\lambda]\in\frakh^*/W$ 的表示的下降数据。特化到奇异的无穷小特征，就恢复出注 4.3 中讨论的从 D-模到表示的商。若再对对称子群引入等变性，则对整个纯内类 $\Theta$ 的 Harish-Chandra 模范畴得到类似的描述
\[
\mathcal{HC}_\Theta\simeq\bigoplus_{\tau\in\Theta}\mathcal D_H(K_\tau\backslash G/N)^{\mathbb D}.
\]
Betti 范畴 $\mathrm{Shv}_N(G/N)$ 与它 de Rham 的对应物 $\mathcal D_H(G/N)$ 的差别，在于丢掉了单值 $\alpha\in\check H$ 的对数 $\lambda\in\frakh^*$ 的选取——这正是环面上局部系统与平坦联络之别：
\[
\Loc(H)\simeq\QCoh(\check H)\qquad\text{与}\qquad \mathrm{Conn}(H)\simeq\QCoh(\frakh^*/X_\bullet(\check H)).
\]
不过，只要记录一个额外的格分级，这一区别就可以消去；特别地，我们可以把 $U\frakg\text{-}\mathrm{mod}$ 恢复为 $\mathrm{Shv}_N(G/N)$ 上某个合适单子的代数。同样地，我们可以恢复
\[
\mathcal{HC}_\Theta\simeq\bigoplus_{\tau\in\Theta}\mathrm{Shv}_N(G_\tau\backslash G/N)^{\mathbb D_{\mathrm{aff}}},
\]
这就把猜想 4.5 的自守侧与实群的表示直接联系了起来。

\subsection{扭量几何 Langlands：谱侧}

Soergel 猜想谱侧的一个自然的族版本，来自考虑扭量线上的 Langlands 参数。我们按如下几种等价的方式引入扭量 Langlands 参数叠 $\mathbb L^\eta$。

(1) 令 $\mathcal B:=\check G/\check B$ 为旗簇，则 $\mathbb L^\eta$ 是叠
\[
\mathbb L^\eta:=\{(g,\check B')\in\check G\times\mathcal B\mid g\eta(g)\in\check B'\}/\check G,
\]
其中 $\check G$ 按 $\eta$-扭共轭作用，即 $h\cdot(g,\check B')=(hg\eta(h^{-1}),h\check B'h^{-1})$。换言之，$\mathbb L^\eta$ 是 $G^L$ 第二个分量上的平方映射与 $\check G/\check G$ 的 Grothendieck--Springer 消解的纤维积\tn{底本图中左下角记作 $\check B/\check B$。按上下文，此处应为映射到 $\check G/\check G$ 的 Grothendieck--Springer 消解的底空间，故该处记号疑有笔误。译文照排。（第 450 页）}：
\[
\begin{tikzcd}
\mathbb L^\eta \arrow[r] \arrow[d] & (G^L-\check G)/\check G \arrow[d,"(-)^2"] \\
\check B/\check B \arrow[r] & \check G/\check G
\end{tikzcd}
\]
存在到泛 Cartan 的单值映射
\[
\chi:\mathbb L^\eta\to\check B/[\check B,\check B]=\check H,\qquad g\mapsto g\eta(g)\pmod{[\check B,\check B]}.
\]

\ysrc{451}

(2) 我们还有用圈空间的描述（见 [BN10b] 的定理 4.6）：
\[
\mathbb L^\eta\simeq \mathcal L((\check B\backslash\check G/\check B)^\Gamma)\simeq \mathcal L(\check B\backslash\check G/\check B)^\Gamma,
\]
其中 $\eta\in\Gamma=\Z/2\Z$ 按 $\eta$ 并交换两个因子而作用；由这一描述，$\mathbb L^\eta$ 带有圆作用。单值映射由 $S^1$-等变映射
\[
\begin{aligned}
\chi:\mathcal L(\check B\backslash\check G/\check B)^\Gamma&\longrightarrow \mathcal L(B\check B\times B\check B)^\Gamma\\
&\simeq\check B/\check B\to\check B//\check B\simeq\check H
\end{aligned}
\]
编码。我们还有 $\eta$-扭圈空间的描述
\[
\begin{aligned}
\mathcal L((\check B\backslash\check G/\check B)^\Gamma)
&\simeq \mathcal L_\eta((\mathcal B\times\mathcal B)/\check G)\\
&=\mathcal L((\mathcal B\times\mathcal B)/G^L)\times_{\mathcal L(B^\Gamma)}\{\eta\},
\end{aligned}
\]
不过这一表述并不先验地从圈空间 $\mathcal L((\mathcal B\times\mathcal B)/G^L)$ 继承 $S^1$-作用（见定义 2.8）\footnote{若施加 $\Gamma$-等变性，即取特殊化的圈空间 $\mathcal L'_\eta$，则有一个圆作用；但这个作用的次数是我们所考虑作用「一半」的次数。$\Gamma$-等变性允许一个良定义的一半次数。}。

(3) 我们还回忆（例如见 [BN16, 3.2]）：$\mathbb L^\eta$ 可以等同于扭量线 $\ol\PP^1_\R$ 上的 $\eta$-扭抛物 $\check G$-局部系统。这里的实形式由 $\PP^1_\C$ 模去由对径映射 $z\mapsto-1/\bar z$ 给出的 $\Gamma$-作用得到。具体地说，我们考虑 $\ol\PP^1_\R\setminus\infty$ 上的 $G^L$-局部系统，其诱导的 $\Gamma$-局部系统等同于定向双覆盖，并且带有包含单值的 Borel 子群。等价地，它们是 $\PP^1_\C\setminus 0,\infty$ 上的 $\check G$-局部系统，在两个极点处有不变的旗，并带有尊重这些旗的 $\Gamma$-等变结构。$\mathbb L^\eta$ 的这个描述使 $S^1$-作用显而易见：它来自 $U(1)$ 作为 $(\ol\PP^1_\R,\infty)$ 的对称群的几何作用（等价地，$\PP^1_\C$ 沿穿过两极点的轴旋转，该旋转与对径映射交换）；事实上，$(\PP^1_\C,0,\infty)$ 上抛物局部系统的叠（即 Steinberg 叠的群版本）被自然地等同于圈空间 $\mathcal L(\check B\backslash\check G/\check B)$。单值映射恰好由绕 $\infty$ 的单值给出。

对泛 Cartan 的元素 $\alpha\in\check H$，定义单值扭量参数空间 $\mathbb L^\eta_\alpha$ 为 $\mathbb L^\eta$ 沿 $\alpha$ 上纤维的形式完备化。[BN10b] 的一个关键观察是：$\mathbb L^\eta_\alpha$ 可等同于相应 ABV 空间的幂幺圈空间（回忆命题 4.1）\footnote{回忆 (4.1) 中 $\Sigma(\eta,\alpha)$ 的含义。}：
\[
\begin{aligned}
\mathbb L^\eta_\alpha&\simeq\bigsqcup_{\sigma\in\Sigma(\eta,\alpha)}\mathcal L^u(K_\sigma(\alpha)\backslash\check G(\alpha)/\check B(\alpha))\\
&\simeq\mathcal L^u((\check B(\alpha)\backslash\check G(\alpha)/\check B(\alpha))^\Gamma).
\end{aligned}
\]
这个等同可视为群 $B\times B$ 的等变局部化的特殊情形（见例 A.2.5）：它由圈空间 $\mathcal L((B\backslash G/B)^\Gamma)$ 在 $\alpha\in B//B=H$ 处完备化而得。特别地，$\mathbb L^\eta_\alpha$ 配备一个 $\GG_m$-作用，把它收缩到 ABV 空间 $X_\alpha$。这个等同在差一个显式中心扭（见第 A.2.3 节）的意义下尊重圆作用，因此循环层的范畴被等同。此外，应用定理 2.7（它把形式圈与幂幺圈上的 Tate 层等同起来），由于 $\GG_m$-作用的收缩性，可以把 $\mathbb L^\eta_\alpha$ 上的循环层直接等同于 ABV 空间上的滤过 D-模（并在有限轨道叠上由 Riemann--Hilbert 得到层），即 Soergel 猜想的谱侧\footnote{关于重新正规化的 Tate 构造的讨论见第 A.1 节。}：

\ysrc{452}

\begin{thm}{4.4}\label{thm:a5-03-4-4}
存在 $k[u,u^{-1}]$-线性的范畴等价
\[
\begin{aligned}
\QCoh^{!}(\mathbb L^\eta_\alpha)^{\mathrm{Tate}}
&\simeq \mathrm{Shv}((\check B(\alpha)\backslash\check G(\alpha)/\check B(\alpha))^\Gamma)\\
&\qquad\otimes_k k[u,u^{-1}],
\end{aligned}
\]
它联系单值扭量 Langlands 参数空间上的周期循环层，与 $\alpha$-ABV 参数空间上的层范畴（基变换到 $k[u,u^{-1}]$）。
\end{thm}

于是我们得到 Soergel 猜想谱侧的一个族版本，可与由 $G/N$ 上幂幺层给出的自守对应物相比较。我们把这一预期编码为下述猜想（暂时推迟 Hecke 作用等自然相容性）：

\begin{conj}{4.5}（族 Soergel 猜想）\label{conj:a5-03-4-5}
存在范畴等价
\[
\mathrm{Shv}_N((N\backslash G/N)^\Gamma)\otimes_k k[u,u^{-1}]\simeq \QCoh^{!}(\mathbb L^\eta)^{\mathrm{Tate}}.
\]
\end{conj}

\subsection{扭量几何 Langlands：自守侧}

猜想 4.2 自守侧出现的叠 $(B\backslash G/B)^\Gamma$（或猜想 4.5 中出现的它的 $N$-版本），在扭量线的意义下有一个自然的几何解释，这由 [BN07] 发现。

\begin{defn}{4.6}\label{defn:a5-03-4-6}
$\ol\PP^1_\R$ 上实 $G$-丛的拓扑叠
\[
\Bun_{G,\theta}(\ol\PP^1_\R;\infty):=\Bun_G(\PP^1;0,\infty)^\Gamma
\]
是如下不动点叠：$\Gamma$ 作用在（$\PP^1$ 上带 0、$\infty$ 处装饰旗（即 $N$-约化）的 $G$-丛叠的 Betti 实现）上，作用由对径映射与 $G$ 上对合 $\theta$ 复合而成。
\end{defn}

几何 Langlands 对应的 Betti 形式 [BN16] 试图用局部系统叠的代数几何来描述 $G$-丛叠上的幂幺层范畴 $\mathrm{Shv}_N(\Bun_G)$。在扭量线上抛物丛的情形，我们已经把相应的局部系统空间与 $\mathbb L^\eta$ 等同起来，于是有下述猜想：

\begin{conj}{4.7}（扭量几何 Langlands [BN07, BN16]）\label{conj:a5-03-4-7}
存在范畴等价
\[
\mathrm{Shv}_N(\Bun_{G,\theta}(\ol\PP^1_\R,\infty))\simeq \QCoh^{!}(\mathbb L^\eta),
\]
它与自然的仿射 Hecke 对称相容。
\end{conj}

\begin{rem}{4.8}\label{rem:a5-03-4-8}
叠上层的大范畴有两种变体（「标准」与「重新正规化」），对应于 $BG$ 上层的两个 Koszul 对偶变体（$G$ 上链代数的模，与 $BG$ 上上链代数的模）。它们分别对应于在谱侧对凝聚层施加或不施加幂零奇异支撑这一条件。
\end{rem}

在 $\Bun_{G,\theta}(\ol\PP^1_\R;\infty)$ 中，底层 $G$-丛在 $\PP^1$ 上平凡的开轨迹，等同于伽罗瓦群在叠 $N\backslash G/N$ 上的不动点叠 $(N\backslash G/N)^\Gamma$；后者参数化平凡丛上成对的装饰旗。因此，Soergel 猜想中出现的范畴 $\mathrm{Shv}_N((N\backslash G/N)^\Gamma)$ 及其诸变体，等同于 $\Bun_{G,\theta}(\ol\PP^1_\R;\infty)$ 某个开轨迹上的层，从而在该范畴的一个\ysrc{453}半正交分解中占一席之地。圆作用尊重这个子范畴，并在其上平凡化；事实上，扭量自守范畴的周期循环形式有如下描述，它来自 [GKM98] 的等变局部化定理，该定理应用于 $\Bun_{G,\theta}(\ol\PP^1_\R,\infty)$ 上 $S^1$-等变层的 $\Hom$ 空间：

\begin{prop}{4.9}（[BN07]）\label{prop:a5-03-4-9}
自守层的周期循环范畴被等同于
\[
\mathrm{Shv}_N(\Bun_{G,\theta}(\ol\PP^1_\R,\infty))^{\mathrm{Tate}}\simeq \mathrm{Shv}_N((N\backslash G/N)^\Gamma)\otimes_k k[u,u^{-1}].
\]
\end{prop}

于是族 Soergel 猜想（猜想 4.5）与扭量几何 Langlands（猜想 4.7）的周期循环形式被等同起来。

换言之，扭量几何 Langlands 猜想（即 $\R$ 上的凝聚局部 Langlands 对应）的 $u$-形变只挑出与 $G$ 的纯内型表示相伴的子范畴，并在该子范畴上产生范畴局部 Langlands 对应的可构造形式。

\begin{rem}{4.10}\label{rem:a5-03-4-10}
对非半单群，平凡丛（或 $S^1$-不动）轨迹——其上的层参与向纯内型表示循环形变的过程——可以严格小于半稳定轨迹，后者的层对应基本等晶体（basic isocrystal）所关联的群的表示。
\end{rem}

\section*{附录 A.\quad 基础}

本节为前几节讨论的表示论应用建立技术基础。我们从一个一般的设定开始：设 $X$ 是特征零域 $k$ 上带仿射对角线的光滑 Artin 叠。

\subsection*{A.1.\quad 圆作用与 $B\GG_a$-作用}

我们对非仿射群叠 $G=S^1,\ B\GG_a,\ B\GG_a\rtimes\GG_m$ 的范畴 $G$-作用作一番离题讨论，因为对习惯仿射代数群作用的读者来说，这些可能颇为陌生。关于此课题的详尽讨论可参看 [Pr11] 第 6 节；$G=B\GG_a$ 的情形我们另请读者参看 [Ch23] 第 3.1 节。我们要提请注意下述两个现象，它们在 $G$ 为仿射代数群时并不出现。

\textbf{（去）等变化对应的缺失：} 由于 $BG=BS^1$ 不是 1-仿射的，$\QCoh(G)$-模范畴 $\mathbf C$ 中并非所有对象都能由不变对象生成，即自然函子
\[
\mathbf C^{\QCoh(G)}\otimes_{\QCoh(BS^1)}\mathbf{Vect}\hookleftarrow\mathbf C
\]
全忠实（命题 A.6），但不再是等价。另一方面，当 $G=B\GG_a$（类似地，$G=B\GG_a\rtimes\GG_m$）时，由 [Ga14] 的定理 2.5.7，$BG=B^2\GG_a$ 是 1-仿射的，故该函子是等价：
\[
\mathbf C^{B\GG_a}\otimes_{\QCoh(B^2\GG_a)}\mathbf{Vect}\xrightarrow{\ \sim\ }\mathbf C.
\]
这一点在第 A.1.1 节讨论。

\ysrc{454}

\textbf{重新正规化的不变量：} 对 $G=S^1,\ B\GG_a$ 与 $B\GG_a\rtimes\GG_m$，叠 $BG$ 不能由完美对象紧生成，即 $\QCoh(BG)\not\simeq\Ind(\Perf(BG))$，并且 ind-完备化与取 $G$-不变量这两个操作不交换。我们将在第 A.1.2 节说明：为什么我们宁愿先取小范畴的 $G$-不变量、再做 ind-完备化（我们称之为紧重新正规化（弱）$G$-不变量范畴），而不是反过来做。

\subsubsection*{A.1.1.\quad 圆作用的线性化与仿射化}

设仿射代数群 $G$ 作用在概形 $X$ 上。对这一设定，可以附着两个不同幺半范畴上的两个模范畴。

\begin{enumerate}[(1)]
\item 范畴 $\QCoh(G)$ 沿群乘法 $m:G\times G\to G$ 的推前是幺半的，并经由沿作用映射 $a:G\times X\to X$ 的推前作用在 $\QCoh(X)$ 上。或者（也更自然），可以把 $\QCoh(G)$ 视为拉回之下的余幺半范畴，并考虑余模范畴。
\item 范畴 $\QCoh(BG)$ 在张量积之下是幺半的，并经由拉回与张量积作用在 $\QCoh(X/G)$ 上。
\end{enumerate}

此外，这两个范畴经由不变量与余不变量的构造相关联：
\[
\QCoh(X)^{\QCoh(G)}\simeq\QCoh(X/G),\qquad \QCoh(X/G)\otimes_{\QCoh(BG)}\mathbf{Vect}_k\simeq\QCoh(X).
\]
这两个操作有时分别称为等变化（equivariantization）与去等变化（de-equivariantization）。

事实上，这些函子可以在一般框架中理解，它们正是 [Ga14] 第 10.2 节中的伴随不变量函子与重构函子：
\[
\begin{tikzcd}
\mathbf{Mod}(\QCoh(BG))\arrow[rr,shift left,"\mathbf D\mapsto\mathbf D\otimes_{\Rep(G)}\mathbf{Vect}_k"]\arrow[rr,shift right,"\mathbf C\mapsto\mathbf C^{\QCoh(G)}"'] & & \mathbf{Comod}(\QCoh(G))
\end{tikzcd}
\]
可以通过典范等价 $\mathbf{Comod}(\QCoh(G))\simeq\mathbf{Mod}(\QCoh(G)^\vee)$ 从较不熟悉的余模过渡到模；若 $\QCoh(G)$ 自对偶，则还有 $\mathbf{Mod}(\QCoh(G)^\vee)\simeq\mathbf{Mod}(\QCoh(G))$；而若 $\QCoh(G)$ 是刚性幺半的，则定义幺半范畴 $\QCoh(G)$ 的结构函子正是定义余幺半范畴的结构函子的右伴随\footnote{在许多情形 $\QCoh(G)$ 都自对偶；例如由 [Ga14] 的 D.1.2，当 $\QCoh(G)^\vee$ 刚性时如此，而由 [GR17] 的命题 3.4.2，任何完美叠都满足这一点。}。我们关心的是这些伴随成为真正等价的条件。按前引文献中的讨论，这在 $BG$ 为 1-仿射时发生；而由 [Ga14] 的命题 10.4.4，当 $G$ 是仿射代数群时 $BG$ 为 1-仿射。这绝非唯一情形；由同前引文献的定理 2.5.7，当 $G=B\GG_a$、$B\GG_a\rtimes\GG_m$ 时 $BG$ 为 1-仿射（且自对偶）。

如前所述，$B^2\GG_a$ 是 1-仿射的，下面将看到 $BS^1$ 不是。我们关心幺半范畴 $\QCoh(S^1)$：其幺半结构由沿乘法 $m:S^1\times S^1\to S^1$ 的推前给出，单位元是恒等元处的摩天层。先回忆下述标准计算。

\ysrc{455}

\begin{prop}{A.1}\label{prop:a5-03-A-1}
范畴 $\QCoh(S^1)$ 等价于 $S^1$ 上取链复形值的局部系统的（导出）范畴。在 Cartier 对偶之下，有幺半等价与等同
\[
\begin{tikzcd}
& \mathbf{Vect} \arrow[dl,"p^*"'] \arrow[d] \arrow[dr,"q_*"] & \\
\QCoh(S^1) \arrow[r,"\sim"] & \mathbf{Mod}(k\Z) \arrow[r,"\sim"] & \QCoh(\GG_m)
\end{tikzcd}
\]
其中 $p:\Spec k\to B\Z=S^1$、$q:\GG_m\to\Spec k$，并且我们在 $\QCoh(S^1)$ 与 $\QCoh(\GG_m)$ 上取张量积作为幺半结构、在 $\mathbf{Mod}(k\Z)$ 上取卷积\footnote{这点对我们并不重要，但请注意 $\QCoh(S^1)^\omega\not\simeq\Perf(S^1)$，故 $S^1$ 不是完美叠。}。特别地，$\QCoh(S^1)$ 自对偶且刚性。
\end{prop}

\begin{rem}{A.2}\label{rem:a5-03-A-2}
为了加深困惑：人们对范畴上的 $S^1$-作用还可能指另一种概念。具体地说，任何拓扑群 $G$ 都可以内化到 $\infty$-范畴中，于是可以用完全抽象的语言表述 $G$ 作用在 $\infty$-范畴上是什么意思\footnote{例如，我们可以定义由单个对象构成的范畴 $BG$，其态射为 $G$，并把 $\mathbf C$ 上的 $G$-作用定义为把该唯一对象送到 $\mathbf C$ 的函子 $BG\to\mathbf{Cat}_\infty$（「拉直」版本）；也可以把 $G$-模范畴定义为余笛卡尔纤维化 $\mathbf C^{\otimes}\to BG$（「未拉直」版本）。}。在 $k$-线性情形，这些概念是等价的。当 $G=S^1$ 时，这样的作用由映射 $\Z\to\HH^\bullet(\mathbf C)$（映入 Hochschild 上同调复形）给出，它等价于 $\mathbf C$ 上的 $\OO(\GG_m)$-线性结构。详见 [Pr11] 第 6.1 节。
\end{rem}

现在我们在等变化的一侧做计算，从 1-仿射的 $BG=B^2\GG_a$ 开始。取 $u\in\OO(B^2\GG_a)$，从而一劳永逸地等同 $k[u]\simeq\OO(B^2\GG_a)$；并令 $\mathbb B=\{0\}\times_{\mathbb A^1}\{0\}$。

\begin{prop}{A.3}\label{prop:a5-03-A-3}
叠 $B^2\GG_a$ 是 1-仿射且自对偶的，因此等变化与去等变化函子都是等价。此外，令 $p:\Spec k\to B^2\GG_a$，有幺半等同
\[
\begin{tikzcd}
& \mathbf{Vect} \arrow[dl,"p^*"'] \arrow[d] \arrow[dr,"R\Hom_{k[u]}(k{,}-)"] & \\
\QCoh(B^2\GG_a) \arrow[r,"\sim"] & \QCoh(\mathbb B) \arrow[r,"\sim"',"-\otimes_{\OO_{\mathbb B}}k"] & \mathbf{Mod}_{u\text{-}\mathrm{tors}}(k[u])
\end{tikzcd}
\]
对 $\QCoh(B^2\GG_a)$ 上的张量幺半结构、$\QCoh(\mathbb B)$ 上的卷积结构，以及 $\mathbf{Mod}_{u\text{-}\mathrm{tors}}(k[u])$（局部 $u$-挠模的满子范畴）上的 $!$-张量积，上述等同成立。特别地，$\QCoh(B^2\GG_a)^\omega\not\simeq\Perf(B^2\GG_a)$，且 $B^2\GG_a$ 不是完美叠\footnote{例如，模 $k[u,u^{-1}]/uk[u]\in\Mod_u(\Sym^\bullet V^*[-2])$ 不紧，但是完美的。}。对群 $G=B\GG_a\rtimes\GG_m$ 有类似的陈述。
\end{prop}

\begin{proof}
为读者方便，我们说明一下 $\QCoh(B^2\GG_a)$ 的计算如何进行（沿 [BN10a]）。首先，由通常的下降论证，$\QCoh(B^2\GG_a)$ 等价于余代数 $\OO(B\GG_a)\simeq\Ext^\bullet_{\GG_a}(k,k)\simeq C^\bullet(S^1;k)$ 的余模，其中余乘法由沿群乘法的拉回给出。对偶之后，这等价于 $C_\bullet(S^1;k)$ 的模。由 Koszul
\ysrc{456}
对偶，这等价于 $C^\bullet(BS^1;k)$ 的局部幂零模（乘法取杯积），即增广模 $k\in\mathbf{Mod}(C^\bullet(BS^1;k))$ 是紧生成元，且 $\End_{C^\bullet(BS^1;k)}(k,k)\simeq C_\bullet(S^1;k)$。
\end{proof}

$BG=BS^1$ 的情形类似（见 [BN10a] 的引理 3.10）；事实上，$B^2\GG_a$ 与 $BS^1$ 上的拟凝聚层范畴没有区别。特别地，$B^2\GG_a$ 是 $BS^1$ 的 1-仿射化。

\begin{prop}{A.4}\label{prop:a5-03-A-4}
仿射化映射 $a:BS^1\to B^2\GG_a$ 诱导幺半等价
\[
a^*:\QCoh(B^2\GG_a)\xrightarrow{\ \sim\ }\QCoh(BS^1).
\]
\end{prop}

经过 Cartier 对偶，在去等变化的一侧我们有幺半范畴的等价
\[
\QCoh(B\GG_a)\simeq\QCoh(\wh{\GG}_a),\qquad \QCoh(S^1)=\QCoh(B\Z)\simeq\QCoh(\GG_m).
\]
这两者显然不同；粗略地说，$\QCoh(S^1)$-模范畴是 $\GG_m$ 上的范畴层，而等变化只能看到 $1\in\GG_m$ 的一个形式邻域。更精确地说，令 $a:BS^1\to B^2\GG_a$ 记仿射化映射\tn{底本此处作 $a:S^1\to B\GG_a$，疑脱去两个 $B$：按命题 A.4，仿射化映射应为 $a:BS^1\to B^2\GG_a$。译文照排。（第 456 页）}、$\iota:\wh{\GG}_a\hookrightarrow\GG_m$，则在 Cartier 对偶之下，伴随对 $(a^*,a_*)$ 与伴随对 $(\iota_*,\iota^!)$ 等同。由于等变化一侧相应的范畴由命题 A.4 等价，我们可以断言：$BS^1$ 不是 1-仿射的，且等变化与去等变化不是互逆等价。

\begin{exam}{A.5}（$BS^1$ 不是 1-仿射的）\label{exam:a5-03-A-5}
考虑正则的 $\QCoh(S^1)$-余模范畴 $\QCoh(S^1)$。它的不变量与增广 $\QCoh(BS^1)$-模范畴 $\mathbf{Vect}_k$ 之间有一个自然的等同，而由命题 A.4 有 $\mathbf{Vect}_k\otimes_{\QCoh(BS^1)}\mathbf{Vect}_k\simeq\QCoh(B\GG_a)$，于是函子
\[
\begin{aligned}
\QCoh(S^1)^{\QCoh(S^1)}\otimes_{\QCoh(BS^1)}\mathbf{Vect}_k
&\simeq\mathbf{Vect}_k\otimes_{\QCoh(BS^1)}\mathbf{Vect}_k\\
&\simeq \QCoh(B\GG_a)\hookrightarrow\QCoh(S^1)
\end{aligned}
\]
全忠实，且由沿仿射化映射 $a:S^1\to B\GG_a$ 的拉回诱导，其本质像是局部幂零 $k\Z$-模构成的满子范畴。特别地，这个函子不是等价，故 $BS^1$ 不是 1-仿射的。
\end{exam}

上述例子中函子的全忠实性在更一般的情形也成立；我们在下一个命题中把这个函子的本质像称为 $S^1$-不变（或 $S^1$-可等变化）对象的满子范畴。注意这个命题有一个重新正规化的对应物，见命题 A.9。

\begin{prop}{A.6}\label{prop:a5-03-A-6}
设 $\mathbf C$ 是 $\QCoh(S^1)$-余模范畴。则
\[
\mathbf C^{\QCoh(S^1)}\otimes_{\QCoh(BS^1)}\mathbf{Vect}\hookrightarrow\mathbf C
\]
全忠实。
\end{prop}

\begin{proof}
右正则表示 $\mathbf C=\QCoh(S^1)$ 的情形已在例 A.5 中确立。一般地，把上述情形与 $\mathbf C$ 作张量：
\[
\mathbf C\otimes_{\QCoh(S^1)}\QCoh(S^1)^{S^1}\otimes_{\QCoh(BS^1)}\mathbf{Vect}\to\mathbf C.
\]
操作 $\mathbf C\otimes_{\QCoh(S^1)}-$ 可通过过渡到右伴随表为极限，因此与 $S^1$-不变量操作交换，并保持全忠实。
\end{proof}
